\section{Partitioned matrices}
A block is itself a matrix. Block addition and multiplication follow the
ordinary row--column rules, but every block product must have matching
inner dimensions. Preserve the order of factors: $AB$ cannot in general
be changed to $BA$. An identity block has the size needed in its own row
and column, so two symbols $I$ in one display may have different sizes.

\practice{2.4.P1}{Invert a block row-replacement matrix}
Let $T=\begin{bmatrix}I&0\\A&I\end{bmatrix}$. Try reversing the addition
of $A$ times the first block row to the second:
\[S=\begin{bmatrix}I&0\\-A&I\end{bmatrix}.\]
Block multiplication gives $TS=\begin{bmatrix}I&0\\A-A&I\end{bmatrix}=I$.
In the reverse order $ST$ has lower-left block $-A+A=0$ and all
other identity blocks. Thus $S=T^{-1}$.

\practice{2.4.P2}{Transpose a horizontal partition before multiplying}
For $X=[X_1\ X_2]$, its transpose is vertical:
\[X^TX=\begin{bmatrix}X_1^T\\X_2^T\end{bmatrix}[X_1\ X_2]
=\begin{bmatrix}X_1^TX_1&X_1^TX_2\\X_2^TX_1&X_2^TX_2\end{bmatrix}.\]
The off-diagonal blocks do not disappear; they measure cross-products
between the two groups of columns.

\exercise{2.4.1}{Add a multiple of one block row}
The conformability assumption makes each displayed product valid.
\[\begin{bmatrix}I&0\\E&I\end{bmatrix}\begin{bmatrix}A&B\\C&D\end{bmatrix}=\begin{bmatrix}A&B\\EA+C&EB+D\end{bmatrix}\]
The lower-left block is $EA+IC=EA+C$; the upper row is unchanged.

\exercise{2.4.2}{Multiply by a block diagonal matrix}
The conformability assumption makes each displayed product valid.
\[\begin{bmatrix}E&0\\0&F\end{bmatrix}\begin{bmatrix}P&Q\\R&S\end{bmatrix}=\begin{bmatrix}EP&EQ\\FR&FS\end{bmatrix}\]
The first block row is multiplied by $E$ and the second by $F$; the zero blocks remove the other terms.

\exercise{2.4.3}{Interchange two block rows}
The conformability assumption makes each displayed product valid.
\[\begin{bmatrix}0&I\\I&0\end{bmatrix}\begin{bmatrix}A&B\\C&D\end{bmatrix}=\begin{bmatrix}C&D\\A&B\end{bmatrix}\]
For example the upper-left product is $0A+IC=C$. This is a block row interchange.

\exercise{2.4.4}{Subtract a block multiple}
The conformability assumption makes each displayed product valid.
\[\begin{bmatrix}I&0\\-E&I\end{bmatrix}\begin{bmatrix}W&X\\Y&Z\end{bmatrix}=\begin{bmatrix}W&X\\Y-EW&Z-EX\end{bmatrix}\]
The minus sign acts on both products in the lower block row. The order $EW$, rather than $WE$, comes from the row--column rule.

\exercise{2.4.5}{Equate the four blocks before choosing inverses}
Multiplication of the two given matrices gives
\[\begin{bmatrix}A+BX&BY\\C&0\end{bmatrix}
=\begin{bmatrix}0&I\\Z&0\end{bmatrix}.\]
Thus the exact requirements are $BX=-A$, $BY=I$, and $Z=C$.
If $B$ is square, these requirements force $B$ invertible, and
\[X=-B^{-1}A,\qquad Y=B^{-1},\qquad Z=C.\]
Indeed $B(-B^{-1}A)=-A$ and $BB^{-1}=I$, checking both upper
blocks. More generally, if $B$ is rectangular and has a right inverse
$Y$, the choice $X=-YA$ still works. Without such a right inverse,
$BY=I$ is impossible, so no solution exists.

\exercise{2.4.6}{Invert a block lower triangular matrix}
The product is
\[\begin{bmatrix}XA&0\\YA+ZB&ZC\end{bmatrix}=I.\]
Assume $A$ and $C$ are square and invertible. Then the diagonal
blocks give $X=A^{-1}$ and $Z=C^{-1}$. The lower-left equation is
$YA=-C^{-1}B$, so right multiplication by $A^{-1}$ gives
\[Y=-C^{-1}BA^{-1}.\]
Inserting these formulas makes $YA+ZB=-C^{-1}B+C^{-1}B=0$.
For square diagonal blocks, their invertibility is also necessary
because $XA=I$ and $ZC=I$ imply it.

\exercise{2.4.7}{A zero middle block does not contribute}
The block product is
\[\begin{bmatrix}XA&XZ\\YA+B&YZ+I\end{bmatrix}=I.\]
With $A$ square and invertible, $XA=I$ gives $X=A^{-1}$.
Then $XZ=0$ forces $Z=0$, and $YA+B=0$ gives $Y=-BA^{-1}$.
The remaining equation is $YZ+I=I$, which is now satisfied.
\answer{$X=A^{-1}$, $Y=-BA^{-1}$, and $Z=0$ under the stated
square-invertible assumption.}

\exercise{2.4.8}{Match three block columns}
Multiplication gives
\[\begin{bmatrix}AX&AY&AZ+B\\0&0&I\end{bmatrix}
=\begin{bmatrix}I&0&0\\0&0&I\end{bmatrix}.\]
If $A$ is square and invertible, the first row requires
\[X=A^{-1},\qquad Y=0,\qquad Z=-A^{-1}B.\]
Each formula follows by multiplying its block equation on the left
by $A^{-1}$. If $A$ is rectangular with a right inverse $X$, one
possible solution is instead $Y=0,Z=-XB$; $AX=I$ checks it.
An $A$ without a right inverse cannot satisfy the first equation.

\exercise{2.4.9}{Eliminate two blocks below a leading block}
The second block row starts with $A_{21}B_{11}+B_{21}$.
Setting it to zero and multiplying on the right by $B_{11}^{-1}$
gives $A_{21}=-B_{21}B_{11}^{-1}$. Likewise
$A_{31}=-B_{31}B_{11}^{-1}$.
The second block in the new second row is consequently
\[C_{22}=A_{21}B_{12}+B_{22}
=B_{22}-B_{21}B_{11}^{-1}B_{12}.\]
For completeness $C_{11}=B_{11}$, $C_{12}=B_{12}$, and
$C_{32}=B_{32}-B_{31}B_{11}^{-1}B_{12}$.
The middle inverse cannot be moved past its neighboring factors.

\exercise{2.4.10}{Invert three block levels in order}
Multiply the original matrix by the proposed inverse. Its lower
blocks must satisfy
\[A+P=0,\qquad D+R=0,\qquad B+DP+Q=0.\]
Thus $P=-A$, $R=-D$, and $Q=DA-B$.
In the reverse product the lower-left block is
$Q+RA+B=(DA-B)-DA+B=0$, while the other off-diagonal blocks
also vanish. Hence both products are the identity.
\answer{$P=-A$, $Q=DA-B$, and $R=-D$.}

\exercise{2.4.11}{Check conformability rather than only the number of blocks}
\begin{enumerate}[label=(\alph*)]
\item \textbf{True.} Corresponding blocks have equal sizes, so adding
entries gives $[A_1+B_1\ A_2+B_2]$.
\item \textbf{False.} The drawn partitions do not specify the actual
sizes. The number of columns of each $A_{ij}$ must equal the number
of rows of the corresponding $B_j$. For example, take all $A_{ij}$
scalar but each $B_j$ a $2\times1$ column. The displayed block
pattern exists but multiplication is not conformable.
\end{enumerate}

\exercise{2.4.12}{Dimensions and transposes of partitions}
\begin{enumerate}[label=(\alph*)]
\item \textbf{False.} $A$ has size $2n\times n$ and $B$ has size
$n\times2n$. Both $BA$ ($n\times n$) and $AB$ ($2n\times2n$)
are defined. The blocks of $AB$ are $A_iB_j$.
\item \textbf{False.} Transposition interchanges block positions as
well as transposing each block:
$A^T=\begin{bmatrix}P^T&R^T\\Q^T&S^T\end{bmatrix}$.
The stated version puts the off-diagonal blocks in the wrong places.
\end{enumerate}

\exercise{2.4.13}{Invertibility of a block diagonal matrix}
If $B$ and $C$ are invertible, block multiplication shows that
$\operatorname{diag}(B^{-1},C^{-1})$ is a two-sided inverse.
Conversely suppose $\operatorname{diag}(B,C)$ is invertible.
If $Bu=0$, then the block vector $(u,0)^T$ is in the kernel of the
whole matrix; injectivity forces $u=0$. Since $B$ is square, the
Invertible Matrix Theorem makes $B$ invertible. The same argument
with $(0,v)^T$ proves $C$ invertible.
\answer{The equivalence holds in both directions.}

\exercise{2.4.14}{Invertibility of a block upper triangular matrix}
Write $A=\begin{bmatrix}A_{11}&A_{12}\\0&A_{22}\end{bmatrix}$ with
square diagonal blocks. If they are invertible, the two-sided inverse is
\[\begin{bmatrix}A_{11}^{-1}&-A_{11}^{-1}A_{12}A_{22}^{-1}\\
0&A_{22}^{-1}\end{bmatrix}.\]
The upper-right block of the product in the original order is
$-A_{12}A_{22}^{-1}+A_{12}A_{22}^{-1}=0$; in the other order
it is $A_{11}^{-1}A_{12}-A_{11}^{-1}A_{12}=0$.
For necessity, invertibility of $A$ first forces $A_{11}$ invertible:
a nonzero $u$ with $A_{11}u=0$ would give kernel vector $(u,0)^T$.
If $A_{22}v=0$ with $v\ne0$, put
$u=-A_{11}^{-1}A_{12}v$. Then $A(u,v)^T=0$, another contradiction.
Thus $A_{22}$ also has trivial kernel and is invertible.

\exercise{2.4.15}{Update an accumulated outer-product sum}
The column--row rule gives
\[G_k=X_kX_k^T=\sum_{j=1}^k x_jx_j^T.\]
Appending the new column yields
\[G_{k+1}=[X_k\ x_{k+1}]\begin{bmatrix}X_k^T\\x_{k+1}^T\end{bmatrix}
=G_k+x_{k+1}x_{k+1}^T.\]
Only one new outer product is needed, followed by entrywise addition
into the stored $G_k$. If each $x_j$ has $m$ entries, the new product
uses $m^2$ scalar products of entries, rather than recomputing $k+1$
outer products. The already processed vectors need not be revisited.

\exercise{2.4.16}{Derive the Schur-complement factorization}
Multiply the last two factors first:
\[\begin{bmatrix}A_{11}&0\\0&S\end{bmatrix}
\begin{bmatrix}I&Y\\0&I\end{bmatrix}
=\begin{bmatrix}A_{11}&A_{11}Y\\0&S\end{bmatrix}.\]
Multiplying the left factor then gives
$\begin{bmatrix}A_{11}&A_{11}Y\\XA_{11}&XA_{11}Y+S\end{bmatrix}$.
Choose $Y=A_{11}^{-1}A_{12}$ and $X=A_{21}A_{11}^{-1}$.
The lower-right block becomes
$A_{21}A_{11}^{-1}A_{12}+S=A_{22}$ by the definition of $S$.
All four blocks now agree with the required matrix.

\exercise{2.4.17}{Invertibility passes to the Schur complement}
The outside factors from Exercise 2.4.16 are invertible: reverse
$X$ or $Y$ in their off-diagonal block, just as in Practice Problem 1.
If $A=L\operatorname{diag}(A_{11},S)U$ is invertible, then
$\operatorname{diag}(A_{11},S)=L^{-1}AU^{-1}$ is invertible.
Exercise 2.4.13 implies $S$ is invertible. Inverting the factors in
reverse order also gives the explicit formula
\[A^{-1}=\begin{bmatrix}
A_{11}^{-1}+A_{11}^{-1}A_{12}S^{-1}A_{21}A_{11}^{-1}&-A_{11}^{-1}A_{12}S^{-1}\\
-S^{-1}A_{21}A_{11}^{-1}&S^{-1}
\end{bmatrix}.\]
This expression follows by multiplying $U^{-1}$, the block diagonal
inverse, and $L^{-1}$, in that order.

\exercise{2.4.18}{Append one data column to a cross-product matrix}
For $W=[X\ x_0]$, transposition and block multiplication give
\[W^TW=\begin{bmatrix}X^TX&X^Tx_0\\x_0^TX&x_0^Tx_0\end{bmatrix}.\]
Its Schur complement of the invertible upper-left block is the scalar
\[x_0^Tx_0-x_0^TX(X^TX)^{-1}X^Tx_0
=x_0^T\bigl[I-X(X^TX)^{-1}X^T\bigr]x_0=x_0^TMx_0.\]
The asserted reciprocal entry requires this scalar to be nonzero;
otherwise $W^TW$ is singular, and an inverse entry is not defined.
The question's requested identity itself holds regardless of whether
the scalar vanishes.

\exercise{2.4.19}{Eliminate the state vector to obtain the transfer function}
Put $H=A-sI_n$, assumed invertible. The top block equation is
$Hx+Bu=0$, so $x=-H^{-1}Bu$. Insert this into the bottom equation
$Cx+u=y$ to obtain
\[\bigl[I_m-CH^{-1}B\bigr]u=y.\]
Thus $W(s)=I_m-C(A-sI_n)^{-1}B$.
It is exactly the Schur complement of $A-sI_n$ in the original
block system matrix. Its size is $m\times m$, so it acts on the
input vector $u\in\R^m$.

\exercise{2.4.20}{The inverse transfer function is another Schur complement}
The Schur complement of $K=A-BC-sI_n$ in the displayed matrix is
\[I_m-(-C)K^{-1}B=I_m+C(A-BC-sI_n)^{-1}B.\]
This is the requested $W(s)^{-1}$ when the stipulated inverses exist.
To check the inverse claim, write $K=H-BC$, $H=A-sI_n$.
From $K=H-BC$, multiply $K K^{-1}B=B$ on the left by
$H^{-1}$ and rearrange to get the identity
\[K^{-1}B=H^{-1}B+H^{-1}BC K^{-1}B.\]
Multiplying it by $C$ shows
$CK^{-1}B=CH^{-1}B+CH^{-1}BC K^{-1}B$.
Expanding $(I-CH^{-1}B)(I+CK^{-1}B)$ now cancels these three
terms and gives $I$, confirming the formula.
The invertibility of $H$ and $W$ also forces $K$ invertible:
if $Kv=0$, then $v=H^{-1}BCv$ and $W(Cv)=0$; hence $Cv=0$
and then $v=0$. The square-matrix theorem completes this check.

\exercise{2.4.21}{Verify an involution using blocks}
\step{(a) Multiply the $2\times2$ matrix.}
\[\begin{bmatrix}1&0\\2&-1\end{bmatrix}^2
=\begin{bmatrix}1&0\\2-2&1\end{bmatrix}=I_2.\]
\step{(b) Recognize the four-dimensional block structure.}
For this same $A$, the given $M$ is
$\begin{bmatrix}A&0\\I_2&-A\end{bmatrix}$. Consequently
\[M^2=\begin{bmatrix}A^2&0\\A-A&A^2\end{bmatrix}=I_4.\]
The lower-left cancellation is essential; squaring only the diagonal
blocks would not verify the entire product.

\exercise{2.4.22}{Construct a six-dimensional involution}
Choose $A=B=I_2$, $D=-I_2$, and
$C=\begin{bmatrix}1&2\\3&4\end{bmatrix}\ne0$. Then
\[M=\begin{bmatrix}I_2&0&0\\0&I_2&0\\C&0&-I_2\end{bmatrix}.
\]
It has size $6\times6$ as required. Its square has identity diagonal
blocks and the only possibly nonzero off-diagonal block is
$CA+DC=C-C=0$. Thus $M^2=I_6$. This construction works with
any nonzero $2\times2$ choice for $C$.

\exercise{2.4.23}{Induct on the dimension of lower triangular matrices}
For dimension one the product of two scalars is lower triangular.
Assume products of $k\times k$ lower triangular matrices are lower
triangular. Partition two matrices of dimension $k+1$ as
\[A_1=\begin{bmatrix}a&0^T\\v&A\end{bmatrix},\qquad
B_1=\begin{bmatrix}b&0^T\\w&B\end{bmatrix}.\]
Their product is
\[A_1B_1=\begin{bmatrix}ab&0^T\\vb+Aw&AB\end{bmatrix}.\]
The top-right block is zero and $AB$ is lower triangular by the
induction hypothesis. These facts account for every entry above the
main diagonal. Thus the product in dimension $k+1$ is lower
triangular, completing the induction.

\exercise{2.4.24}{Invert cumulative sums by successive differences}
Let $A_n$ be the lower triangular matrix of ones and $B_n$ the
matrix with ones on the diagonal and $-1$ immediately below it.
For $n=2$, direct multiplication of
$\begin{bmatrix}1&0\\1&1\end{bmatrix}$ and
$\begin{bmatrix}1&0\\-1&1\end{bmatrix}$ in either order gives $I_2$.
Suppose $A_kB_k=B_kA_k=I_k$. With $\mathbf1$ the $k$-entry
ones vector and $e_1=(1,0,\ldots,0)^T$, partition
\[A_{k+1}=\begin{bmatrix}1&0^T\\\mathbf1&A_k\end{bmatrix},\quad
B_{k+1}=\begin{bmatrix}1&0^T\\-e_1&B_k\end{bmatrix}.\]
The lower-left block of $A_{k+1}B_{k+1}$ is
$\mathbf1-A_ke_1=0$, since $A_ke_1$ is the first column, all ones.
The lower-left block of the reverse product is
$-e_1+B_k\mathbf1=0$: successive differences of ones are
$(1,0,\ldots,0)^T$. The diagonal blocks are identities by induction
and the upper-right blocks are zero. Both products are $I_{k+1}$,
which proves the assertion for every $n\ge2$.

\exercise{2.4.25}{Invert the three diagonal blocks separately}
The given matrix is block diagonal with blocks
$B=\begin{bmatrix}1&2\\3&5\end{bmatrix}$, $[2]$, and
$C=\begin{bmatrix}7&8\\5&6\end{bmatrix}$.
The $2\times2$ formula gives
\[B^{-1}=\begin{bmatrix}-5&2\\3&-1\end{bmatrix},\qquad
C^{-1}=\begin{bmatrix}3&-4\\-5/2&7/2\end{bmatrix}.\]
Thus, without row reduction,
\[A^{-1}=\begin{bmatrix}-5&2&0&0&0\\3&-1&0&0&0\\
0&0&1/2&0&0\\0&0&0&3&-4\\0&0&0&-5/2&7/2\end{bmatrix}.\]
Each product block is an identity, so the whole product is $I_5$.

\exercise{2.4.26}{Concrete commands for extracting and inserting blocks}
Using Python/NumPy's zero-based, end-exclusive slices:
\begin{enumerate}[label=(\alph*)]
\item \verb|A[4:10, 14:20]| displays rows 5--10 and columns 15--20.
The result has size $6\times6$.
\item \verb|A[4:9, 9:19] = B| inserts the $5\times10$ block in rows
5--9 and columns 10--19. The first inserted entry is at row 5,
column 10, as required.
\item \verb|scipy.linalg.block_diag(A, A.T)| constructs the desired
$50\times50$ block diagonal matrix. Alternatively initialize a
$50\times50$ zero array, assign $A$ to \verb|C[:20, :30]| and
$A^T$ to \verb|C[20:, 30:]|. The off-diagonal blocks then already
have the necessary zero entries.
\end{enumerate}
These are the requested operations, not a request for unspecified data.
Their output sizes follow from the slice endpoints and the given sizes.

\exercise{2.4.27}{Work entirely with blocks below the size limit}
Keep the large matrices on disk or as separately stored blocks;
assembling a full $50\times50$ array would itself violate the limit.
\begin{enumerate}[label=(\alph*)]
\item Partition both matrices into four $25\times25$ blocks.
Compute each sum block $C_{ij}=A_{ij}+B_{ij}$ separately and
store the four results as the sum's block representation.
\item With that same partition compute
$C_{ij}=A_{i1}B_{1j}+A_{i2}B_{2j}$ for $i,j\in\{1,2\}$.
Every individual multiplication and addition is $25\times25$,
well within the $32\times32$ bound.
\item Use the stipulated $20+30$ partition, with
$b=(b_1,b_2)^T$ and $x=(x_1,x_2)^T$. Since $A_{12}=0$,
first solve the $20\times20$ system $A_{11}x_1=b_1$.
Compute $r=b_2-A_{21}x_1$ ($A_{21}$ is $30\times20$), then
solve $A_{22}x_2=r$, a $30\times30$ system. Both diagonal
matrices are invertible, so both solves have unique answers.
Store $x_1,x_2$ separately as the solution; no explicit inverse
or overlarge matrix is needed.
\end{enumerate}
